\chapter{Máquinas e Instalaciones Eléctricas: Controladores de Potencia y Accionamientos}
\label{sec:maquinas_e_instalaciones_electricas}

En la asignatura Máquinas e Instalaciones Eléctricas, el anteproyecto desarrolla el diseño electromecánico y de
potencia para los accionamientos de la CNC. Esto abarca los controladores modulares de velocidad y sentido de giro para
los motores paso a paso de los tres ejes y el regulador de potencia para la herramienta de corte (motor DC).

\section{Requerimiento de Diseño Modular por Canal}
\label{sec:myie_modularidad}

Para garantizar un mantenimiento sencillo, una escalabilidad flexible y cumplir con la consigna académica, el sistema
de potencia adopta una topología estrictamente modular.

Cada eje mecánico (X, Y, Z) cuenta con su propio módulo electrónico independiente montado en una tarjeta PCB canalizada.
Cada canal de potencia integra:
\begin{itemize}
  \item La etapa lógica de decodificación de las señales digitales STEP y DIR.
  \item Las etapas de potencia compuestas por puentes H para la conmutación de las fases del motor paso a paso.
  \item La resistencia de derivación ($R_{shunt\_canal}$) colocada en el retorno de masa común del canal, a fin de permitir
    que la etapa analógica de Electrónica Aplicada II tome la medición de corriente total consumida por dicho canal.
\end{itemize}

En la figura~\ref{crkt.ej4:driver_paso_a_paso} se presenta la estructura de bloques de un canal modular de potencia.

\begin{figure}[!ht]
  \centering
  \begin{tikzpicture}[
  >=stealth, thick,
  font=\small,
  node distance=1.2cm and 2.0cm,
  block/.style={draw=orange!80!black, fill=orange!12, rectangle, rounded corners=4pt, minimum height=1.2cm, minimum width=2.9cm, align=center},
  ctrl/.style={draw=blue!70!black, fill=blue!10, rectangle, rounded corners=4pt, minimum height=1.0cm, minimum width=2.5cm, align=center},
  sens/.style={draw=purple!70!black, fill=purple!10, rectangle, rounded corners=4pt, minimum height=0.9cm, minimum width=2.9cm, align=center},
  motor/.style={draw=green!60!black, fill=green!10, rectangle, rounded corners=4pt, minimum height=1.2cm, minimum width=2.6cm, align=center, font=\small\bfseries}
]

  % Bloques principales
  \node[ctrl] (ctrl) {Señales de Control\\{\footnotesize (STEP / DIR / EN)}};
  \node[block, right=1.6cm of ctrl] (logic) {Lógica de Control\\y Modulación\\{\footnotesize (Secuenciador / Chopper)}};
  \node[block, right=2.0cm of logic] (power) {Etapa de Potencia\\{\footnotesize (Puentes H / Fases A y B)}};
  \node[sens, below=1.2cm of power] (sens) {Retorno de Masa y\\Punto de Sensado};
  \node[motor, right=2.0cm of power] (motor) {Motor PaP\\{\footnotesize (Eje X / Y / Z)}};

  % Contenedor Modular (Caja con margen superior amplio de 1.8cm para no colisionar con la alimentación)
  \draw[dashed, draw=gray!60, rounded corners=6pt] 
    ([xshift=-0.5cm, yshift=2.0cm]logic.north west) rectangle 
    ([xshift=0.5cm, yshift=-0.4cm]sens.south east);
  \node[anchor=north west, font=\footnotesize\bfseries\color{black!75}] 
    at ([xshift=-0.35cm, yshift=1.85cm]logic.north west) {Módulo de Potencia por Canal (PCB Independiente)};

  % Conexiones con etiquetas centradas
  \draw[->] (ctrl) -- (logic);
  \draw[->] (logic) -- node[above, font=\footnotesize] {Comando} (power);
  \draw[->] (power) -- node[above, font=\footnotesize] {Fases} (motor);
  \draw[->] (power) -- (sens);
  \draw[->, draw=purple!80!black] (sens.east) -| node[pos=0.25, below, font=\footnotesize] {Hacia Etapa EA2} ([xshift=-0.3cm]motor.south);

  % Alimentación con espacio holgado
  \draw[<-] (power.north) -- ++(0,0.8) node[above, font=\footnotesize] {Alimentación $V_{mot}$};

\end{tikzpicture}

  \caption{Diagrama de bloques funcional de un canal modular de potencia para motor paso a paso.}
  \label{crkt.ej4:driver_paso_a_paso}
\end{figure}

\section{Controlador de Motores Paso a Paso y Pasos Fraccionados (Micropasos)}
\label{sec:myie_paso_a_paso}

Los motores paso a paso recuperados de las impresoras suelen ser del tipo bipolar de dos fases. El controlador del
canal genera las secuencias de activación necesarias mediante las señales de entrada:
\begin{itemize}
  \item \textbf{Señal STEP:} Cada flanco de subida de esta señal determina el avance de un paso (o subpaso) angular del
    motor. La frecuencia de la señal impone directamente la velocidad de rotación ($n$ en $\mathrm{rpm}$).
  \item \textbf{Señal DIR:} Define la secuencia de fase aplicada al motor, determinando si el giro se realiza en sentido
    horario o antihorario.
\end{itemize}

\subsection{Evaluación de Complejidad para Pasos Fraccionados}
Un aspecto central de estudio en la materia es el análisis de implementar modulación por troceado de corriente
(\textit{chopping PWM}) para lograr pasos fraccionados (micropasos de 1/2, 1/4 o 1/8). 

En un accionamiento por pasos completos, la conmutación abrupta de tensión genera vibraciones y resonancia en la
estructura liviana de la impresora. La implementación de micropasos aproxima la corriente en las fases a ondas senoidales
desfasadas $90^\circ$, tal como se expresa en la ecuación~\ref{eq.ej4:fases_micro}:

\begin{equation}
  I_A(\theta) = I_{max} \cdot \sin(\theta), \quad I_B(\theta) = I_{max} \cdot \cos(\theta)
  \label{eq.ej4:fases_micro}
\end{equation}

El anteproyecto plantea la evaluación de complejidad para determinar si se utilizará control analógico por chopper
dedicado o generación de consigna sinusoidal desde el microcontrolador, ponderando el compromiso entre suavidad de
movimiento, torque disponible y tiempo de desarrollo.

\section{Controlador del Motor DC para el Spindle (Fresa)}
\label{sec:myie_spindle}

El cabezal de mecanizado utiliza un motor de corriente continua de alta velocidad para impulsar la fresa o mecha. La
regulación de su velocidad resulta indispensable para adaptar las condiciones de corte al material trabajado.

El circuito de potencia utiliza un MOSFET de potencia actuando como conmutador de estado sólido en serie con la carga,
controlado por una señal PWM generada desde la placa digital. La figura~\ref{crkt.ej4:driver_spindle} muestra el
esquema circuital de la etapa de potencia del motor DC.

\begin{figure}[!ht]
  \centering
  \begin{tikzpicture}[
  >=stealth, thick,
  font=\small,
  node distance=1.0cm and 2.2cm,
  ctrl/.style={draw=blue!70!black, fill=blue!10, rectangle, rounded corners=4pt, minimum height=1.1cm, minimum width=2.6cm, align=center},
  block/.style={draw=orange!80!black, fill=orange!12, rectangle, rounded corners=4pt, minimum height=1.1cm, minimum width=2.8cm, align=center},
  motor/.style={draw=green!60!black, fill=green!10, rectangle, rounded corners=4pt, minimum height=1.1cm, minimum width=2.6cm, align=center, font=\small\bfseries}
]

  \node[ctrl] (ctrl) {Consigna PWM\\{\footnotesize (Temporizador MCU)}};
  \node[block, right=2.2cm of ctrl] (driver) {Acondicionamiento /\\Driver de Compuerta};
  \node[block, right=2.2cm of driver] (switch) {Conmutación y\\Supresión Inductiva};
  \node[motor, right=2.8cm of switch] (motor) {Motor DC\\{\footnotesize (Spindle / Fresa)}};

  % Conexiones lineales con espacio amplio para las etiquetas
  \draw[->] (ctrl) -- node[above, font=\footnotesize] {PWM} (driver);
  \draw[->] (driver) -- node[above, font=\footnotesize] {Disparo} (switch);
  \draw[->] (switch) -- node[above, font=\footnotesize] {Tensión Media} (motor);

  % Alimentación de potencia
  \draw[<-] (switch.north) -- ++(0,0.8) node[above, font=\footnotesize] {Alimentación $V_{DC}$};

\end{tikzpicture}

  \caption{Esquema circuital de accionamiento de potencia PWM para el motor DC de la herramienta de corte.}
  \label{crkt.ej4:driver_spindle}
\end{figure}

El diodo en paralelo inverso con la inductancia del motor (\textit{flyback diode}) protege al transistor de sobretensiones
inductivas provocadas por el apagado rápido del MOSFET durante la conmutación PWM. La velocidad del motor responde de forma
directamente proporcional al ciclo de trabajo ($\delta$) de la señal de control.
